\section{The actual first controlled stage in the original units}
\label{sec:cv-first-stage}

The source for this section is Alp\"oge--Buckmaster,
\emph{Blowup for the Boussinesq equations with smooth forcing},
\url{https://cims.nyu.edu/~tristanb/boussinesq.pdf}, released PDF,
76 pages. The source locations used here are Lemma 1.2 and (1.6),
pp.~6--7; (3.1)--(3.12), pp.~8--10; the complete proofs of
Lemmas 3.7--3.8 and (3.24)--(3.34), pp.~15--18;
(3.58) and the following second-stage calculation, p.~27; and
the first holding integral on p.~36. The proofs below are given in
full. The source's unit-scale convention is not imposed on the fields
in this section.

\subsection{Parameters, coordinates, and the original coupled equations}

Let $x=(x_1,x_2)\in\mathbb R^2$, $e_2=(0,1)$, and
\[
 J=\begin{pmatrix}0&-1\\1&0\end{pmatrix},\quad
 R_\alpha=\begin{pmatrix}\cos\alpha&-\sin\alpha\\
 \sin\alpha&\cos\alpha\end{pmatrix},\quad
 e(s)=(\sin s,\cos s).
\]
In particular $R_\alpha e(s)=e(s-\alpha)$ and
$\operatorname{curl}u=\partial_1u_2-\partial_2u_1$.
Retain $A_0>0$, $\lambda_0\geq1$, and the source's two scaling parameters
\[
 \mu=\lambda_0,\qquad \nu=\sqrt{A_0},\qquad K_1=\mu\lambda_1.
\]
Here $\nu$ is the source's \emph{time-scaling parameter}.
Physical momentum viscosity below is denoted $\nu_{\rm phys}\geq0$;
thermal diffusivity is $\kappa_{\rm phys}\geq0$.
The first pulse compression is the source's $\Lambda_1$, which is
different from the physical spatial frequency $K_1$.

Fix a smooth nondecreasing step $\eta:\mathbb R\to[0,1]$ equal to zero
on $(-\infty,0]$ and one on $[1,\infty)$, and
$h\in C_c^\infty((0,1))$ with $h\geq0$ and $\int_0^1h=1$.
Write $H=\|h\|_\infty$. The source's choice is $h=\eta_2$,
$c_p=3$; the following calculation retains any
\begin{equation}
 c_p>1,\quad \Lambda_1\geq2c_pH,\quad
 0<s=s_*\leq\tfrac18,\quad c_p\Lambda_1H\sin s\leq\tfrac14.
 \label{eq:cv-first-design}
\end{equation}
Since $\int h=1$, $H\geq1$ and $\Lambda_1>2$.
Let $k_1\geq0$ be an integer, $\lambda_1>1$, and set
\begin{equation}
 S=\frac{A_0}{\mu}\lambda_1^{-k_1-6},\quad
 L=(k_1+5+\tfrac18)\log\lambda_1\geq1,\quad
 P_* =S e^L=\frac{A_0}{\mu}\lambda_1^{-7/8},\quad
 \Gamma=\nu\sin s,\quad t_0=\nu^{-1}.
 \label{eq:cv-first-seed}
\end{equation}
The condition $L\geq1$ is ensured explicitly by
$\lambda_1\geq\exp((k_1+41/8)^{-1})$.

The first layer's complete preceding affine background, phase transport,
and signed amplitude pair are
\begin{align}
 D_{<1}&=\dot\alpha J,&G_{<1}&=-A_0R_\alpha e_2,
 &\dot M_1&=D_{<1}M_1,&M_1(t_0)&=I,\label{eq:cv-first-background}\\
 \zeta_1&=M_1^{-T}e(s),&
 \dot\Theta_1&=a_1\Omega_1,&
 \dot\Omega_1&=b_1\Theta_1,&
 a_1&=-\frac{J\zeta_1\cdot G_{<1}}{K_1|\zeta_1|^2},\quad
 b_1=K_1\zeta_{1,1}.\label{eq:cv-first-pair}
\end{align}
The entry data are
\[
 \alpha(t_0)=0,\quad \Theta_1(t_0)=-S,\quad
 \Omega_1(t_0)=-\frac{K_1}{\nu}S,
 \qquad w_1(t)=\eta(\Gamma(t-t_0)).
\]
The activation factor multiplies the physical increment, not the
homogeneous amplitude equation. It is flat at $t_0$ and equals one
by $t_0+\Gamma^{-1}$.

For this actual background, $M_1=R_\alpha$: differentiation gives
$\dot R_\alpha=\dot\alpha JR_\alpha$ and the initial values agree.
Uniqueness follows, for example, by differentiating
$R_\alpha^{-1}M_1$. Consequently
\begin{equation}
 \zeta_1=R_\alpha e(s),\quad |\zeta_1|=1,\quad
 J\zeta_1\cdot G_{<1}=-A_0\sin s,\quad
 a_1=\frac{A_0\sin s}{K_1},\quad
 b_1=K_1\sin(s-\alpha).
 \label{eq:cv-first-exact-coefficients}
\end{equation}
The determinant identity used here is explicit:
$J e(s)=(-\cos s,\sin s)$ has scalar product $\sin s$
with $e_2$, and simultaneous rotation preserves that scalar product.
In particular $a_1$ is constant because both the parent gradient and
the new phase rotate; it has not been frozen independently of the
evolving parent.

\subsection{Exact growth and transition}

For the first stage the feedback sums $F_0$ and $F_1$ in source (3.4)
are empty and equal zero. Thus both the source's growth control
$\dot\alpha=F_0$ and its transition control
$(1-\eta_b)F_0+\eta_bF_1$ give $\alpha=0$.
Set
\[
 t_a=t_0+\frac L\Gamma,\qquad
 t_1=t_a+\Gamma^{-1},\qquad
 t_b=t_1+\Gamma^{-1}(1+\Lambda_1^{-1}).
\]
On $[t_0,t_1]$ the exact solution of the original pair is
\begin{equation}
 P_1(t):=-\Theta_1(t)=S e^{\Gamma(t-t_0)},\qquad
 \Omega_1(t)=-\frac{K_1}{\nu}P_1(t).
 \label{eq:cv-first-growth}
\end{equation}
Indeed $a_1K_1/\nu=\Gamma$ and
$\Gamma K_1/\nu=K_1\sin s=b_1$.
This proves both differential equations and both initial data.
Uniqueness follows from the linear integral equation: on a bounded
interval with coefficient norm at most $B$, the $n$-fold iterated
integral of the difference of two solutions is bounded by its supremum
times $(B\ell)^n/n!$, which tends to zero. The same convergent integral
series supplies existence on every finite interval.
Since $P_1$ is strictly increasing, $t_a$ is exactly the first crossing
of $P_*$, and the transition contributes exactly one to $\log P_1$.
The steering entry amplitude is therefore $P_{\rm ent}=eP_*$.

\subsection{The trial steering family and proof of its unique return}

Use $m\in[0,c_p]$ for the pulse trial parameter, to distinguish it
from the unchanged spatial parameter $\mu=\lambda_0$.
With $\tau=\Gamma(t-t_1)$ and $\tau_b=1+\Lambda_1^{-1}$ set
\begin{equation}
 z(\tau;m)=
 \begin{cases}
  1-\eta(\tau),&0\leq\tau\leq1,\\
  -m\Lambda_1h(\Lambda_1(\tau-1)),&1\leq\tau\leq\tau_b,
 \end{cases}
 \quad Z=\sin s\,z,\quad \alpha=s-\arcsin Z.
 \label{eq:cv-first-Z}
\end{equation}
Extend $Z=\sin s$ before $t_1$ and $Z=0$ after $t_b$.
The two smooth pieces and extensions agree to every order, since the
step is flat and $h$ is supported in the interior.
The source's selected pulse height in (3.12) is precisely
$h_1=m\Lambda_1\sin s$; its allowed interval is
$[0,c_p\Lambda_1\sin s]$.
The design gives
\begin{equation}
 |Z|\leq\tfrac14,\quad
 \zeta_{1,1}=Z,\quad
 \zeta_{1,2}=\sqrt{1-Z^2}\geq\frac{\sqrt{15}}4,\quad
 \dot\alpha=-\frac{\dot Z}{\zeta_{1,2}},\quad 0\leq\alpha<\tfrac12.
 \label{eq:cv-first-control}
\end{equation}
For the last bound, $Z\leq\sin s$ implies $\alpha\geq0$, and
$\arcsin(1/4)\leq(1/4)/\sqrt{1-1/16}<1/3$ gives
$\alpha\leq1/8+1/3<1/2$. These are the source's actual feedback and
the original laboratory phase component; gravity remains $e_2$.

Define $\mathcal P$ first as the solution of the linear equation
\begin{equation}
 \mathcal P_{\tau\tau}=z\mathcal P,\qquad
 \mathcal P(0)=\mathcal P_\tau(0)=1.
 \label{eq:cv-first-linear-multiplier}
\end{equation}
The same linear integral series proves existence, uniqueness, and
smoothness for every trial on the full finite interval. Moreover the
exact inverse map to the original amplitude pair is
\begin{equation}
 \Theta_1(t)=-P_{\rm ent}\mathcal P(\tau),\qquad
 \Omega_1(t)=-\frac{K_1}{\nu}P_{\rm ent}\mathcal P_\tau(\tau).
 \label{eq:cv-first-inverse-map}
\end{equation}
To check it, its component derivatives are exactly
\[
 -\Gamma P_{\rm ent}\mathcal P_\tau=a_1\Omega_1,\qquad
 -\frac{K_1}{\nu}\Gamma P_{\rm ent}z\mathcal P
 =K_1\sin s\,z\Theta_1=b_1\Theta_1.
\]
The data at $t_1$ match \eqref{eq:cv-first-growth}. Conversely the
first amplitude equation with constant $a_1$ and $a_1b_1=\Gamma^2z$
gives \eqref{eq:cv-first-linear-multiplier} for
$-\Theta_1/P_{\rm ent}$, with the stated initial derivative.
Thus the auxiliary equation and original pair are explicitly inverse
descriptions of the same dynamics, not a substituted oscillator.

On $0\leq\tau\leq1$, $z\geq0$. Before any first zero of either
$\mathcal P$ or $\mathcal P_\tau$, one has
$\mathcal P_{\tau\tau}\geq0$ and therefore
$\mathcal P_\tau\geq1$, $\mathcal P\geq1+\tau$.
Continuity excludes that first zero and proves these inequalities
through $\tau=1$. Only now define $v=\mathcal P_\tau/\mathcal P$.
It satisfies
\[
 v_\tau=1-\eta-v^2,\qquad v(0)=1.
\]
Since $(v-1)_\tau=-(v+1)(v-1)-\eta$, the integrating-factor formula
gives $v\leq1$. Since $v>0$ and $v_\tau\geq-v^2$,
$(1/v)_\tau\leq1$, so
\begin{equation}
 \frac1{1+\tau}\leq v(\tau)\leq1,\quad
 \log2\leq\int_0^1v\,d\tau\leq1,\quad
 v_d:=v(1)\in[1/2,1],\quad 2\leq\mathcal P(1)\leq e.
 \label{eq:cv-first-ramp}
\end{equation}

On the negative pulse set $y=\Lambda_1(\tau-1)$ and
$\chi(y)=\mathcal P(1+y/\Lambda_1)/\mathcal P(1)$.
There is no division by an unknown evolving amplitude in this
definition. Its exact equation and data are
\begin{equation}
 \chi_{yy}+\frac{mh(y)}{\Lambda_1}\chi=0,\qquad
 \chi(0)=1,\quad\chi_y(0)=v_d/\Lambda_1.
 \label{eq:cv-first-chi}
\end{equation}
Before a possible first zero, $\chi$ is concave, so
$\chi\leq1+y/\Lambda_1\leq2$. Since
$mh/\Lambda_1\leq c_pH/\Lambda_1\leq1/2$, one has
$\chi_{yy}\geq-1$, hence
\[
 \chi(y)\geq1+\frac{v_d}{\Lambda_1}y-\frac{y^2}{2}\geq\frac12
 \qquad(0\leq y\leq1).
\]
The strict positive lower bound excludes a first zero. Thus every
trial satisfies on the entire pulse
\begin{equation}
 \tfrac12\leq\chi(y)\leq1+y/\Lambda_1,\qquad
 \chi_y=\frac{v_d}{\Lambda_1}
  -\frac m{\Lambda_1}\int_0^yh(r)\chi(r)\,dr
  \geq-\frac{2c_p}{\Lambda_1}.
 \label{eq:cv-first-trial-positive}
\end{equation}
Consequently the original temperature is strictly negative for every
trial, and the quotient is well defined. Writing
$V(y;m)=v(1+y/\Lambda_1;m)=\Lambda_1\chi_y/\chi$ gives
\begin{equation}
 V_y=-mh(y)-\frac{V^2}{\Lambda_1},\quad V(0;m)=v_d,
 \qquad -4c_p\leq V(y;m)\leq v_d\leq1.
 \label{eq:cv-first-exact-riccati}
\end{equation}

Here is the parameter dependence needed to select the return.
In the first-order integral system for $(\chi,\chi_y)$ the coefficient
matrix depends affinely on $m$. On the compact parameter interval,
its iterated integrals converge uniformly with factorial denominators.
Differentiating $r$ times in $m$ introduces at most a degree-$r$
polynomial in the iteration index, still dominated by that factorial.
Thus this series and each parameter derivative converge uniformly.
The lower bound $\chi\geq1/2$ then gives smooth parameter dependence
of $V$. Differentiating its exact equation and using an integrating
factor gives
\begin{equation}
 \partial_m V(1;m)=
 -\int_0^1h(y)\exp\!\left(-\frac2{\Lambda_1}
                 \int_y^1V(r;m)\,dr\right)dy<0.
 \label{eq:cv-first-root-derivative}
\end{equation}
The strict sign follows because $h\geq0$ has mass one and the
exponential is positive. Indeed the explicit bounds also give
$e^{-2/\Lambda_1}\leq-\partial_mV(1;m)
\leq e^{8c_p/\Lambda_1}$.
At $m=0$, direct solution of $V_y=-V^2/\Lambda_1$ gives
$V(1;0)=v_d/(1+v_d/\Lambda_1)>0$. At $m=c_p$, integration gives
$V(1;c_p)\leq v_d-c_p<0$. Continuity and strict monotonicity
therefore give exactly one $m_*\in(0,c_p)$ with $V(1;m_*)=0$.
This is exactly the source's selected return: by
\eqref{eq:cv-first-inverse-map} and positivity, it is equivalent
to $\Omega_1(t_b)=0$.

For this selected trial $V$ is nonincreasing and ends at zero, so
$0\leq V\leq v_d$. Integration of
\eqref{eq:cv-first-exact-riccati} now gives the exact identity and bounds
\begin{equation}
 m_*=v_d-\frac1{\Lambda_1}\int_0^1V(y;m_*)^2\,dy,
 \qquad \frac12-\frac1{4\Lambda_1}\leq m_*\leq1.
 \label{eq:cv-first-root-bounds}
\end{equation}
For the lower bound, $v_d-v_d^2/\Lambda_1$ is increasing for
$v_d\in[1/2,1]$, because its derivative is
$1-2v_d/\Lambda_1>0$.

For every trial and every steering time,
\[
 0\leq\log\mathcal P(\tau)\leq1+\Lambda_1^{-1}.
\]
On the ramp this follows from \eqref{eq:cv-first-ramp}; on the pulse
it follows from $\mathcal P(1)\geq2$, $\chi\geq1/2$, and
$\log\chi\leq\log(1+\Lambda_1^{-1})\leq\Lambda_1^{-1}$.
For the selected trial the pulse logarithmic contribution is
nonnegative, so the endpoint obeys
\begin{equation}
 \log2\leq\log\mathcal P(\tau_b)\leq1+\Lambda_1^{-1},\qquad
 2eP_*\leq P_b:=P_1(t_b)\leq e^{2+\Lambda_1^{-1}}P_*.
 \label{eq:cv-first-gain}
\end{equation}
The total duration from entry to return is exactly
\begin{equation}
 t_b-t_0=\frac{L+2+\Lambda_1^{-1}}{\nu\sin s}.
 \label{eq:cv-first-duration}
\end{equation}
Throughout the selected stage $P_1$ is nondecreasing,
$S\leq P_1\leq P_b$, and
$-(K_1/\nu)P_1\leq\Omega_1\leq0$.
For every unselected trial the valid bound instead is
$|\Omega_1|\leq4c_p(K_1/\nu)
 e^{2+\Lambda_1^{-1}}P_*$; its sign need not stay negative.

Near $t_b$, $h=0$, so $Z=0$ and $\alpha=s$.
The equation $\mathcal P_{\tau\tau}=0$, together with
$\mathcal P_\tau(\tau_b)=0$, implies that $\mathcal P$ is constant
on that whole neighborhood. Thus $\alpha$, $\Theta_1$, and $\Omega_1$
match their stationary continuations to all orders. The endpoint state is
\begin{equation}
 \alpha=s,\quad M_1=R_s,\quad\zeta_1=e_2,\quad
 w_1=1,\quad\Theta_1=-P_b,\quad\Omega_1=0.
 \label{eq:cv-first-endpoint}
\end{equation}

\subsection{Control bounds and exact whole-stage integrals}

For $j=1,2$ define the explicit finite profile quantities
\[
 B_j=\Gamma^j\sin s\bigl(\|\eta^{(j)}\|_\infty
            +c_p\Lambda_1^{j+1}\|h^{(j)}\|_\infty\bigr).
\]
From \eqref{eq:cv-first-Z}, $|\partial_t^jZ|\leq B_j$.
Differentiating the original control, including its denominator, gives
\begin{align}
 \dot\alpha&=-\dot Z(1-Z^2)^{-1/2},\\
 \ddot\alpha&=-\ddot Z(1-Z^2)^{-1/2}
              -Z\dot Z^2(1-Z^2)^{-3/2},\nonumber\\
 |\dot\alpha|&\leq\frac4{\sqrt{15}}B_1,\qquad
 |\ddot\alpha|\leq\frac4{\sqrt{15}}B_2
                    +\frac{16}{15\sqrt{15}}B_1^2.
 \label{eq:cv-first-control-bounds}
\end{align}
Higher derivatives exist because the displayed smooth composition has
argument in $[-1/4,1/4]$; no endpoint differentiability is omitted.

For exact coefficient costs use the invertible fixed map
$Y=(\Theta_1,\nu\Omega_1/K_1)^T$. Its inverse is
$(Y_1,K_1Y_2/\nu)^T$, and direct differentiation gives
\[
 \dot Y=\widehat B Y,\qquad
 \widehat B=\begin{pmatrix}0&\Gamma\\\nu Z&0\end{pmatrix}.
\]
Let $J_\eta=\int_0^1(1-\eta(r))\,dr\in[0,1]$.
On a hold of physical duration $T_h\geq0$, the following are exact:
\begin{center}
\begin{tabular}{c|c|c|c}
interval&duration&$\int|\widehat B_{12}|\,dt$
                      &$\int|\widehat B_{21}|\,dt$\\\hline
growth&$L/\Gamma$&$L$&$L$\\
transition&$1/\Gamma$&$1$&$1$\\
steering ramp&$1/\Gamma$&$1$&$J_\eta$\\
negative pulse&$1/(\Lambda_1\Gamma)$&$1/\Lambda_1$&$m_*$\\
holding&$T_h$&$\Gamma T_h$&$0$
\end{tabular}
\end{center}
For example, on the pulse substitution
$y=\Lambda_1\Gamma(t-t_1-\Gamma^{-1})$ gives
$\int\nu|Z|\,dt=m_*\int_0^1h=m_*$.
Thus the two transition-and-steering costs are exactly
$2+\Lambda_1^{-1}$ and $1+J_\eta+m_*$;
their sum is at most $4+c_p+\Lambda_1^{-1}$.

The activation and amplitude integrals needed by force estimates can
also be retained without an asymptotic replacement. Set
\[
 J_a=\int_0^1\eta'(r)e^r\,dr\in[1,e],\qquad
 J_g=\int_0^1\eta(r)e^r\,dr=e-J_a.
\]
For the selected stage including the hold, direct change of variable
and integration by parts give
\begin{align}
 \int_{t_0}^{t_b}|\dot w_1\Theta_1|\,dt&=S J_a,\qquad
 \int_{t_0}^{t_b}|\dot w_1\Omega_1|\,dt=\frac{K_1}{\nu}S J_a,
 \label{eq:cv-first-activation-cost}\\
 I_P:=\int_{t_0}^{t_b+T_h}w_1P_1\,dt
 &=\frac S\Gamma(J_g+e^{L+1}-e)+P_bT_h\nonumber\\
 &\quad+\frac{eP_*}{\Gamma}\int_0^{1+\Lambda_1^{-1}}
                    \mathcal P(\tau;m_*)\,d\tau,
 \label{eq:cv-first-P-cost}\\
 I_\Omega:=\int_{t_0}^{t_b+T_h}w_1|\Omega_1|\,dt
 &=\frac{P_b-SJ_a}{a_1}.
 \label{eq:cv-first-Omega-cost}
\end{align}
For the last identity, $P_1'=-a_1\Omega_1=a_1|\Omega_1|$,
$w_1(t_0)=0$, and $w_1(t_b+T_h)=1$, hence
$\int w_1P_1'=P_b-\int\dot w_1P_1=P_b-SJ_a$.
In particular
\[
 I_P\leq\frac S\Gamma(J_g+e^{L+1}-e)
 +\frac{eP_*}{\Gamma}(1+\Lambda_1^{-1})e^{1+\Lambda_1^{-1}}
 +P_bT_h.
\]
These are finite specified integrals of the prescribed profiles and
the uniquely selected exact linear solution.

\subsection{The evolving background inherited by layer two}

On the affine part of the first wave, the complete background for the
next layer is
\begin{equation}
 G_{<2}=-A_0R_\alpha e_2+w_1K_1\Theta_1\zeta_1,
 \qquad
 D_{<2}=\dot\alpha J+w_1\Omega_1J\zeta_1\otimes\zeta_1.
 \label{eq:cv-first-inherited}
\end{equation}
These are source (3.3) with $A_1=-K_1\Theta_1|\zeta_1|$
and the physical units retained. Their complete derivatives are
\begin{align}
 \dot G_{<2}&=\dot\alpha JG_{<2}
      +(K_1\dot w_1\Theta_1+w_1A_0\sin s\,\Omega_1)\zeta_1,
 \label{eq:cv-first-G-derivative}\\
 \dot D_{<2}&=\ddot\alpha J
  +(\dot w_1\Omega_1+w_1K_1\zeta_{1,1}\Theta_1)
                            J\zeta_1\otimes\zeta_1\nonumber\\
 &\hspace{5mm}+w_1\Omega_1\dot\alpha
       (-\zeta_1\otimes\zeta_1+J\zeta_1\otimes J\zeta_1).
 \label{eq:cv-first-D-derivative}
\end{align}
Indeed $\dot\zeta_1=\dot\alpha J\zeta_1$,
$J\dot\zeta_1=-\dot\alpha\zeta_1$, and the product rule gives
each displayed term. In particular, using
$J\zeta_1\cdot G_{<2}=-A_0\sin s$ gives the exact identity
\begin{equation}
 \dot G_{<2}+D_{<2}^TG_{<2}=K_1\dot w_1\Theta_1\zeta_1.
 \label{eq:cv-first-affine-scalar-residual}
\end{equation}
Thus the inherited gradient satisfies its true affine transport
equation after activation, even while the first stage rotates it.
The affine vorticity is $2\dot\alpha+w_1\Omega_1$ and its scalar
source is $G_{<2,1}=A_0\sin\alpha+w_1K_1\Theta_1\zeta_{1,1}$;
therefore its exact vorticity residual is
\begin{equation}
 (2\dot\alpha+w_1\Omega_1)'-G_{<2,1}
       =2\ddot\alpha-A_0\sin\alpha+\dot w_1\Omega_1.
 \label{eq:cv-first-affine-vorticity-residual}
\end{equation}

At the selected endpoint define $A_1=K_1P_b>0$. Then
\begin{equation}
 G_{<2}(t_b)=(A_0\sin s,-A_0\cos s-A_1),\quad D_{<2}(t_b)=0,
 \quad\sigma_1^2=|G_{<2}(t_b)|
   =\sqrt{A_0^2+A_1^2+2A_0A_1\cos s}.
 \label{eq:cv-first-exact-sigma}
\end{equation}
Consequently
\[
 2e A_0\lambda_1^{1/8}\leq A_1
 \leq e^{2+\Lambda_1^{-1}}A_0\lambda_1^{1/8},\qquad
 A_1+A_0\cos s\leq\sigma_1^2\leq A_1+A_0.
\]
The lower norm bound is its positive vertical component in absolute
value; the upper bound is the triangle inequality, which also follows
by squaring the displayed expression. No horizontal component of the
parent gradient has been discarded.

On the first hold the actual feedback is $F_1=0$.
The endpoint state \eqref{eq:cv-first-endpoint} is stationary under
the original equations: $b_1=K_1\zeta_{1,1}=0$ gives
$\dot\Omega_1=0$, then $\dot\Theta_1=a_1\Omega_1=0$,
and $\dot\alpha=0$ gives $\dot\zeta_1=0$.
Thus \eqref{eq:cv-first-exact-sigma} remains the actual inherited
background throughout this hold. For an arbitrary perturbation of
the weighted amplitude pair its propagator over length $T_h$ is
exactly
\begin{equation}
 \exp\!\left(T_h\begin{pmatrix}0&\Gamma\\0&0\end{pmatrix}\right)
 =\begin{pmatrix}1&\Gamma T_h\\0&1\end{pmatrix}.
 \label{eq:cv-first-hold-propagator}
\end{equation}
This follows by squaring the matrix and obtaining zero, and retains
the nonzero first holding coefficient cost.

For clarity one can construct the next growth interval itself rather
than prescribing an unspecified holding time. Retain
$\lambda_2=\lambda_1^{Q_2}$ with $Q_2\geq202$, an integer $k_2\geq0$,
$L_2=(k_2+41/8)\log\lambda_2$, and the source's
\[
 s_2=L_2\nu/\sigma_1,\qquad K_2=\mu\lambda_2,
 \qquad\Gamma_2=\sigma_1\sin s_2.
\]
An explicit sufficient finite choice ensuring $L_2\geq2$ and
$0<s_2\leq1/8$ is
$\log\lambda_1\geq4096(k_2+41/8)Q_2$.
To prove this, put $B=(k_2+41/8)Q_2$ and $x=\log\lambda_1$.
The gain implies $\sigma_1\geq\sqrt{2e}\nu e^{x/16}$;
the exponential series gives $e^{x/16}\geq x^2/512$.
Hence $s_2\leq512B/(\sqrt{2e}x)\leq1/8$.
Also $L_2=Bx\geq4096B^2>2$.

During the second growth, $\alpha=s$, $\Omega_1=0$, and
$D_{<2}=0$, so $M_2=I$, $\zeta_2=e(s_2)$, and the original
coefficient formulas yield exactly
\begin{equation}
 a_2=\frac{A_1\sin s_2+A_0\sin(s+s_2)}{K_2},\quad
 b_2=K_2\sin s_2,\quad
 \gamma_2=\sqrt{[A_1\sin s_2+A_0\sin(s+s_2)]\sin s_2}>0.
 \label{eq:cv-first-second-growth}
\end{equation}
Both signs are positive because $0<s+s_2\leq1/4$.
For its actual seed $S_2=(A_0/\mu)\lambda_2^{-k_2-6}$ and
entry vorticity $-\sqrt{b_2/a_2}\,S_2$, the exact solution is
$P_2=S_2e^{\gamma_2(t-t_b)}$,
$\Omega_2=-\sqrt{b_2/a_2}\,P_2$.
Substitution verifies both equations; the first threshold time is
$t_{a,2}=t_b+L_2/\gamma_2$.
Its activation is the source's $\eta(\Gamma_2(t-t_b))$.
To see that activation completes before threshold, note
\begin{equation}
 R:=\frac{\gamma_2^2}{\Gamma_2^2}
  =\frac{A_1+A_0\cos s+A_0\sin s\cot s_2}{\sigma_1^2},
 \qquad
 \cos s\leq R\leq1+\frac{\nu\sin s}{L_2\sigma_1}.
 \label{eq:cv-first-rate-ratio}
\end{equation}
For the lower bound, discard the positive cotangent term, use
$\sigma_1^2\leq A_1+A_0$, and then
$(A_1+A_0\cos s)/(A_1+A_0)\geq\cos s$.
For the upper bound use $A_1+A_0\cos s\leq\sigma_1^2$ and
$\cot s_2\leq1/s_2$, the latter following from
$(\sin r-r\cos r)'=r\sin r\geq0$.
Since $\sigma_1\geq\nu$ and $L_2\geq2$,
$R\leq1+\sin s/2<2$, and
$\Gamma_2L_2/\gamma_2=L_2/\sqrt R>1$.
The exact first holding integral therefore is
\begin{equation}
 H_1=\Gamma\frac{L_2}{\gamma_2}
     =\sin s\,\frac{s_2}{\sin s_2}\frac1{\sqrt R},\quad
 \frac{\sin s}{\sqrt{1+\nu\sin s/(L_2\sigma_1)}}
 \leq H_1\leq
 \frac{\sin s}{\sqrt{\cos s}(1-s_2^2/6)}.
 \label{eq:cv-first-exact-holding-cost}
\end{equation}
Here $\sin r\leq r$ follows by integrating $\cos r\leq1$;
$\sin r\geq r-r^3/6$ follows by integrating
$\cos r\geq1-r^2/2$, itself obtained from
$1-\cos r=\int_0^r\sin u\,du\leq r^2/2$.
Thus no small-angle replacement is used in the bounds.
The second transition also has $F_2=0$, because its only shear
coefficient is $\Omega_1=0$; it prolongs the first pair's stationarity
by exactly $\Gamma_2^{-1}$. The source counts that transition in
the first layer's tail, not in $H_1$.

\subsection{Exact uncut forced fields and their retained diffusion cost}

This subsection gives an explicit PDE realization of the just-proved
finite stage on $\mathbb R^2$. Its fields have the spatial domain
$\mathbb R^2$ and need not decay; compact spatial localization is a
separate construction. Let $F$ be any smooth odd $2\pi$-periodic
function. Define $Q_0(r)=\int_0^rF(v)\,dv$ and
$\mathcal Q(r)=Q_0(r)-(2\pi)^{-1}\int_{-\pi}^{\pi}Q_0(v)\,dv$.
Oddness gives zero period integral for $F$, so $Q_0$ is periodic;
substitution $v=-u$ gives $Q_0(-r)=Q_0(r)$. Thus $\mathcal Q$ is
even and mean-zero, and $\mathcal Q'=F$.
Put $q(x,t)=K_1\zeta_1(t)\cdot x$ and
\begin{align}
 \theta&=G_{<1}\cdot x+w_1\Theta_1F(q),\nonumber\\
 u&=\dot\alpha Jx+\frac{w_1\Omega_1}{K_1}J\zeta_1F(q),
 \label{eq:cv-first-uncut-fields}\\
 p&=\frac{\dot\alpha^2}2|x|^2
      +\frac12x_2(G_{<1}\cdot x)
      +\left(\frac{2w_1\Omega_1\dot\alpha}{K_1^2}
             +\frac{w_1\Theta_1\zeta_{1,2}}{K_1}\right)\mathcal Q(q).
 \label{eq:cv-first-uncut-pressure}
\end{align}
An additive function of time can be subtracted if $p(0,t)=0$ is desired.
The exact inviscid forces are
\begin{equation}
 f_\theta^0=\dot w_1\Theta_1F(q),\qquad
 f_u^0=\left(\ddot\alpha-\frac{A_0}2\sin\alpha\right)Jx
                   +\frac{\dot w_1\Omega_1}{K_1}J\zeta_1F(q).
 \label{eq:cv-first-uncut-forces}
\end{equation}
We verify every component. The background derivative satisfies
$\dot G_{<1}=\dot\alpha JG_{<1}$, hence its scalar transport
residual vanishes. The phase satisfies
$(\partial_t+\dot\alpha Jx\cdot\nabla)q=0$.
The wave velocity has scalar product zero with $\nabla q$, so it
advects neither its own temperature nor its own velocity profile.
Its advection of the background temperature equals
$(w_1\Omega_1/K_1)(J\zeta_1\cdot G_{<1})F(q)$,
which cancels $w_1\dot\Theta_1F(q)$ by
\eqref{eq:cv-first-pair}. The remaining scalar residual is exactly
$\dot w_1\Theta_1F(q)$.

For momentum, write $C=w_1\Omega_1/K_1$ and
$v=CJ\zeta_1F(q)$. Differentiating $v$ and adding its transport
of $\dot\alpha Jx$ gives, after subtracting the wave buoyancy,
\[
 \left[\frac{\dot w_1\Omega_1}{K_1}J\zeta_1
   +\frac{w_1\dot\Omega_1}{K_1}J\zeta_1
   -2C\dot\alpha\zeta_1-w_1\Theta_1e_2\right]F(q).
\]
The last three vector terms have zero scalar product with $J\zeta_1$:
this is $w_1\dot\Omega_1/K_1-w_1\Theta_1\zeta_{1,1}=0$.
Their scalar product with $\zeta_1$ is
$-2C\dot\alpha-w_1\Theta_1\zeta_{1,2}$.
Since $(\zeta_1,J\zeta_1)$ is an orthonormal basis, the wave pressure
in \eqref{eq:cv-first-uncut-pressure} cancels precisely these three
terms, leaving the displayed activation force. The base momentum
residual before pressure is
$(\ddot\alpha J-\dot\alpha^2I-e_2\otimes G_{<1})x$.
The gradient of the first two pressure terms removes its symmetric
part. Its antisymmetric part is
$(\ddot\alpha-G_{<1,1}/2)Jx
=(\ddot\alpha-A_0\sin\alpha/2)Jx$, as claimed.
Finally $\operatorname{div}u=0$ since $\operatorname{tr}J=0$
and $J\zeta_1\cdot\zeta_1=0$.

For the retained-viscosity equations
\[
 \partial_t\theta+u\cdot\nabla\theta
     =\kappa_{\rm phys}\Delta\theta+f_\theta,\qquad
 \partial_tu+(u\cdot\nabla)u+\nabla p
     =\nu_{\rm phys}\Delta u+\theta e_2+f_u,
 \qquad\nabla\cdot u=0,
\]
the very same state and controls solve them with the explicitly
specified compensating forces
\begin{align}
 f_\theta&=f_\theta^0
 -\kappa_{\rm phys}w_1\Theta_1K_1^2F''(q),\nonumber\\
 f_u&=f_u^0-\nu_{\rm phys}w_1\Omega_1K_1J\zeta_1F''(q).
 \label{eq:cv-first-diffusion-compensation}
\end{align}
Indeed the affine Laplacians vanish, while two spatial derivatives
of the profile multiply by $K_1^2|\zeta_1|^2=K_1^2$.
This proves the asserted equations by substitution, without changing
either $\nu$ or $\nu_{\rm phys}$ and without assuming damping leaves
the parent coefficients unchanged. The exact cumulative $L^\infty_x$
norms of the two added forces through the entire stage and hold are
\begin{equation}
 \kappa_{\rm phys}K_1^2\|F''\|_\infty I_P,
 \qquad
 \nu_{\rm phys}K_1\|F''\|_\infty I_\Omega.
 \label{eq:cv-first-exact-diffusion-cost}
\end{equation}
Equality holds because $x\mapsto K_1\zeta_1\cdot x$ is surjective
onto $\mathbb R$ and $|J\zeta_1|=1$; this verifies the norm statement
as well as an upper bound. The scalar holding contribution is exactly
$\kappa_{\rm phys}K_1^2\|F''\|_\infty P_bT_h$; the velocity holding
contribution is zero. These finite costs are not asserted to be
summable over the source's infinite frequency sequence.

The construction proves a finite controlled stage, its actual evolving
affine background, its return, and its next growth holding interval.
It does not establish an infinite viscous cascade or a smooth
terminal-time viscous force. In particular, simply replacing the
first temperature by an exponentially damped amplitude during holding
would replace $A_1$ in \eqref{eq:cv-first-exact-sigma} by a
time-dependent quantity and would change both coefficients and the
rate in \eqref{eq:cv-first-second-growth}. The compensating forces
in \eqref{eq:cv-first-diffusion-compensation} give an exact specified
relation that retains the original coupled finite trajectory; their
costs have been retained explicitly rather than claimed negligible.
